@book{fowler2002patterns,
  title={Patterns of Enterprise Application Architecture},
  author={Fowler, Martin},
  year={2002},
  publisher={Addison-Wesley Longman},
  address={Boston}
}

@article{sinnott1984virtues,
  title={Virtues of the Haversine},
  author={Sinnott, Roger W},
  journal={Sky and Telescope},
  volume={68},
  number={2},
  pages={159},
  year={1984}
}

@phdthesis{fielding2000architectural,
  title={Architectural styles and the design of network-based software architectures},
  author={Fielding, Roy Thomas},
  year={2000},
  school={University of California, Irvine}
}

@book{gurnani2020developing,
  title={Developing Applications with React Native},
  author={Gurnani, Akshat},
  year={2020},
  publisher={Packt Publishing Ltd},
  address={Birmingham}
}

@book{sommerville2011software,
  title={Software Engineering},
  author={Sommerville, Ian},
  year={2011},
  publisher={Addison-Wesley},
  address={Boston}
}

% ==============================================================================
% MANUAL DE OPERAÇÃO: GERENCIAMENTO DE REFERÊNCIAS (BIBTEX + ABNTEX2)
% ==============================================================================
% 
% 1. O CONCEITO DO ARQUIVO .BIB
% No LaTeX, você não digita a lista de referências manualmente no final do texto.
% O arquivo 'referencias.bib' atua como um banco de dados relacional. Você
% cadastra os metadados das fontes nele e faz "consultas" no texto principal.
% Regra de Ouro: Uma fonte cadastrada no .bib SÓ aparecerá no PDF final se
% for explicitamente citada em algum lugar do texto.
%
% 2. COMO CADASTRAR UMA FONTE (No arquivo referencias.bib)
% Cada entrada possui um "tipo" (@book, @article, @misc) e uma "chave" única
% que você inventa para chamá-la depois.
%
% Exemplo de Livro Técnico:
% @book{sommerville2011,
%   author    = {Sommerville, Ian},
%   title     = {Engenharia de software},
%   publisher = {Pearson Prentice Hall},
%   year      = {2011},
%   address   = {São Paulo}
% }
%
% Exemplo de Site, Documentação ou Repositório:
% @misc{fastapi2026,
%   author        = {Ramírez, Sebastián},
%   title         = {FastAPI Documentation},
%   year          = {2026},
%   url           = {https://fastapi.tiangolo.com/},
%   urlaccessdate = {8 jun. 2026}
% }
% Nota: Autores corporativos (ex: Meta, Google, W3C) devem ser envolvidos 
% em chaves duplas no arquivo .bib para que o LaTeX não inverta os nomes. 
% Ex: author = {{Meta Platforms}}
%
% 3. COMO CHAMAR A FONTE NO TEXTO (Nos arquivos .tex)
% O pacote abntex2cite fornece comandos rigorosos para adequação à NBR 10520.
%
% -> Citação Indireta (Ao final da frase, em maiúsculas):
% Comando: \cite{chave}
% Exemplo: A arquitetura baseada em eventos reduz o acoplamento \cite{sommerville2011}.
% PDF renderizado: A arquitetura baseada em eventos reduz o acoplamento (SOMMERVILLE, 2011).
%
% -> Citação Direta/Narrativa (Inserida no fluxo do texto):
% Comando: \citeonline{chave}
% Exemplo: Segundo \citeonline{sommerville2011}, a engenharia de requisitos é vital.
% PDF renderizado: Segundo Sommerville (2011), a engenharia de requisitos é vital.
%
% -> Inserção Forçada (Sem citação textual):
% Comando: \nocite{chave}
% Uso: Insere a fonte na bibliografia final sem precisar colocar um (AUTOR, Ano) no texto.
% Evite usar isso em trabalhos acadêmicos para burlar a fundamentação.
% ==============================================================================